\chapter{Appendix}

\section{Auxiliary Lemmas}\label{app:lemmas}

\begin{lemma}[Parity selector via contour integral]\label{lem:kronecker-even-delta}
For $n\in\mathbb{N}_0$ and any $0<r<1$, define
\[
I_n \;=\; \frac{1}{2\pi i}\oint_{|z|=r}\frac{dz}{z^{\,n+1}(1-z^2)}.
\]
Then
\[
I_n \;=\; \begin{cases}
1,& \text{if $n$ is even},\\[4pt]
0,& \text{if $n$ is odd}.
\end{cases}
\]
\begin{proof}
Let $h(z)=\dfrac{1}{1-z^2}$. For $|z|<1$,
\[
h(z)=\sum_{m=0}^{\infty} z^{2m}.
\]
Writing $h(z)=\sum_{k=0}^{\infty} a_k z^k$, we have $a_k=1$ if $k$ is even and $a_k=0$ if $k$ is odd. By Cauchy’s coefficient formula,
\[
I_n=\frac{1}{2\pi i}\oint_{|z|=r}\frac{h(z)}{z^{n+1}}\,dz=a_n,
\]
so $I_n=1$ when $n$ is even and $I_n=0$ when $n$ is odd.

\medskip
\noindent\textit{Residue viewpoint.}
The integrand $\dfrac{1}{z^{n+1}(1-z^2)}$ has singularities at $z=0,\pm1$, but for $0<r<1$ only $z=0$ lies inside the contour. Hence the integral collapses to the single residue at $z=0$:
\[
I_n \;=\; \operatorname{Res}_{z=0}\!\left(\frac{1}{z^{n+1}(1-z^2)}\right).
\]
Expanding $\dfrac{1}{1-z^2}=\sum_{m=0}^{\infty}z^{2m}$ gives the Laurent series
\[
\frac{1}{z^{n+1}(1-z^2)}=\sum_{m=0}^{\infty} z^{\,2m-n-1},
\]
so the residue (coefficient of $z^{-1}$) is $1$ iff $2m=n$, and $0$ otherwise. Equivalently, using $\operatorname{Res}_{z=0}\!\big(\frac{h(z)}{z^{n+1}}\big)=\frac{h^{(n)}(0)}{n!}$ for analytic $h$ with $h(z)=\frac{1}{1-z^2}$, we have $h^{(n)}(0)=n!$ for even $n$ and $0$ for odd $n$, yielding the same conclusion.
\end{proof}
\end{lemma}

\begin{lemma}[Odd parity selector]\label{rem:odd-selector}
By a similar construction, we can derive a kernel that extracts \emph{odd} coefficients. Define:
\[
J_n \;=\; \frac{1}{2\pi i}\oint_{|z|=r}\frac{z\,dz}{z^{\,n+1}(1-z^2)} = \frac{1}{2\pi i}\oint_{|z|=r}\frac{dz}{z^{n}(1-z^2)}.
\]
Let $g(z) = \dfrac{z}{1-z^2}$. For $|z|<1$,
\[
g(z) = z \sum_{m=0}^{\infty} z^{2m} = \sum_{m=0}^{\infty} z^{2m+1}.
\]
Writing $g(z)=\sum_{k=0}^{\infty} b_k z^k$, we have $b_k=1$ if $k$ is odd and $b_k=0$ if $k$ is even. Therefore:
\[
J_n = \frac{1}{2\pi i}\oint_{|z|=r}\frac{g(z)}{z^{n+1}}\,dz = b_n = \begin{cases}
1,& \text{if $n$ is odd},\\[4pt]
0,& \text{if $n$ is even}.
\end{cases}
\]

Thus the \emph{odd selection kernel} is:
\[
K_N^{\text{odd}}(\zeta) = \sum_{n=0}^{N-1} \zeta^{2n+1} = \zeta \cdot \frac{1 - \zeta^{2N}}{1 - \zeta^2} = \frac{\zeta(1-\zeta^{2N})}{1-\zeta^2},
\]
which selects coefficients $\{a_1, a_3, a_5, \ldots, a_{2N-1}\}$ when used in the projector integral.
\end{lemma}

\begin{lemma}[$\mathcal{H}_S(z)$ does not allow any further elementary decomposition]

\begin{proof}
It is known that any 3 ${}_1F_2$ series with the same argument, with parameters differing by integers, are linearly related\cite{}

Let $x \;:=\; -\frac{z^2}{4}.$
\begin{equation}\label{eq:H_raise}
x\,\mathcal{H}_{S+2}(z)
=\left(\frac{S}{2}+1\right)\left(\frac{S}{2}+\frac{3}{2}\right)\big[\mathcal{H}_S(z)-1\big].
\tag{R+}
\end{equation}

\begin{equation}\label{eq:H_lower}
x\,\mathcal{H}_S(z)
=\left(\frac{S}{2}\right)\left(\frac{S}{2}+\frac{1}{2}\right)\big[\mathcal{H}_{S-2}(z)-1\big].
\tag{R-}
\end{equation}

\begin{equation}\label{eq:H_contiguous}
\left(\frac{S}{2}+1\right)\left(\frac{S}{2}+\frac{3}{2}\right)\mathcal{H}_{S+2}(z)
-\left(\frac{S}{2}\right)\left(\frac{S}{2}+\frac{1}{2}\right)\mathcal{H}_{S-2}(z)
= -\left[\left(\frac{S}{2}\right)\left(\frac{S}{2}+\frac{1}{2}\right)+x\right].
\tag{C}
\end{equation}

\[
\Big[\phi(\phi+b-1)(\phi+c-1)-x(\phi+1)\Big]\,y=0,
\qquad
\phi:=x\frac{d}{dx},\quad b=\frac{S}{2}+1,\quad c=\frac{S}{2}+\frac{3}{2}.
\]
\end{proof}
\end{lemma}

\begin{lemma}[Hypergeometric Form for Real Arguments]\label{app:hypergeometric}

The hypergeometric representation of the truncation remainder derived in \ref{}.
\[
\mathbb{P}_N[\mathrm{sinc}](z) = \mathrm{sinc}(z) + \frac{(-1)^N z^{2(N+1)}}{\Gamma(2N+4)} \mathcal{H}_{2(N+1)}(z)
\]
with
\[
\mathcal{H}_S(z) = {}_1F_2\!\left(1;\,\frac{S}{2}+1,\;\frac{S}{2}+\frac{3}{2};\,-\frac{z^2}{4}\right)
\]
is valid \emph{only for $z \in \mathbb{R}$}.

\begin{proof}[Proof of restriction] In the derivation, we used:
\[
\Im\!\left(\theta^{S-1} \sum_{k=1}^{\infty} \frac{\theta^{k}}{(S)_k}\right) = -\,i\,\theta^{S-1}\sum_{k=1}^{\infty} \frac{\theta^{2k}}{(S)_{2k}}, \quad \theta = iz, \; S = 2(N+1).
\]
This step relies on the fact that $\theta^{S-1} = (iz)^{2N+1}$ is purely imaginary (since $2N+1$ is odd), allowing us to conclude that only even powers of $\theta$ contribute to $\Im(\cdot)$.

However, for $z \in \mathbb{C} \setminus \mathbb{R}$, we have $\theta = iz \notin i\mathbb{R}$, so $\theta^{2k}$ is \emph{not} purely real, and the separation:
\[
\Im\!\left(\sum_{k=1}^{\infty} \frac{\theta^{k}}{(S)_k}\right) = \text{(only even $k$ terms)}
\]
\emph{fails}. The correct approach for complex $z$ is to retain the full complex form:
\[
\mathbb{P}_N[\mathrm{sinc}](z) = \frac{\Im\!\big(e^{iz}\Gamma(2N+2,iz)\big)}{z\,(2N+1)!}, \quad z \in \mathbb{C}.
\]
\end{proof}
\end{lemma}

\section{Design of the general support kernel $\hat{H}(\omega)$}\label{app:support-kernel}

\subsection{Example of shaping kernel using standard design method: digital First-Order Low-Pass Filter}\label{app:first-order-lpf}

Consider the analog first-order low-pass filter with impulse response $h_a(t)$. Its Laplace transform is:
\[
\mathcal{L}\{h_a(t)\}(s) = H_a(s) = \frac{\omega_c}{s + \omega_c}, \quad s \in \mathbb{C}.
\]

Applying the bilinear transformation $s \mapsto \frac{2}{T}\frac{1-z^{-1}}{1+z^{-1}}$ with sampling period $T$ yields the z-transform:
\[
H(z) = \frac{\omega_c T(1+z^{-1})}{(2 + \omega_c T) + (-2 + \omega_c T)z^{-1}}.
\]

Let $\Omega := \omega_c T$ denote the normalized cutoff frequency. Dividing numerator and denominator by $(2 + \Omega)$:
\[
H(z) = \frac{\frac{\Omega}{2 + \Omega}(1+z^{-1})}{1 + \frac{-2 + \Omega}{2 + \Omega}z^{-1}}.
\]

From $\beta := \frac{\Omega}{2 + \Omega}$, we obtain $\beta(2 + \Omega) = \Omega$, hence $\Omega = 2\beta + \beta\Omega$. Rearranging:
\[
\Omega(1 - \beta) = 2\beta \quad \Longrightarrow \quad \Omega = \frac{2\beta}{1-\beta} = 2\beta \sum_{k=0}^{\infty} \beta^k, \quad |\beta| < 1.
\]

Therefore:
\[
\frac{-2 + \Omega}{2 + \Omega} = -(1-\beta) + \beta = 2\beta - 1,
\]
yielding:
\begin{equation}
H(z) = \frac{\beta(1+z^{-1})}{1 - (1-2\beta)z^{-1}}
= \frac{\beta(z + 1)}{2\beta + z - 1}, \quad \text{with pole } z_p = 1 - 2\beta
\label{eq:digital_lpf}
\end{equation}

The convergence condition $|\beta| < 1$ from the geometric series ensures $|z_p| < 1$ (stability), demonstrating equivalence between algebraic convergence and z-domain stability.

\subsection{Special decomposition cases for $\hat{H}_a(\omega)$}\label{app:special-decomp}

We consider the expression:
\[
I = \oint_C \frac{1}{2} [(S_e - iz S_o)H_a(\omega) + (S_e + iz S_o)H_a(-\omega)] \, d\omega
\]
where $S_e$ is a purely even series, $S_o$ is a purely odd series, $\omega = iz/\zeta$ (change of variables), and $C$ is a contour in the complex $\omega$-plane.

\subsubsection{Case 1: Rational Functions}\label{app:case-rational}

Consider $H_a(\omega) = \frac{P(\omega)}{Q(\omega)} = \sum_{k=1}^{N} \frac{R_k}{\omega - p_k}$ where $p_k$ are the poles and $R_k$ are the residues from partial fraction decomposition. Then:
\begin{align*}
I &= \frac{1}{2}\oint_C \left[(S_e - izS_o)\sum_{k=1}^{N}\frac{R_k}{\omega - p_k} + (S_e + izS_o)\sum_{k=1}^{N}\frac{R_k}{-\omega - p_k}\right] d\omega \\
&= \frac{1}{2}\sum_{k=1}^{N} R_k \oint_C \left[\frac{S_e - izS_o}{\omega - p_k} + \frac{S_e + izS_o}{-\omega - p_k}\right] d\omega
\end{align*}

For symmetric contours and poles in pairs $\pm p_k$ with equal residues:
\[
\boxed{I = \pi i \sum_{k: p_k \in C} R_k \left[2S_e + iz S_o \cdot \text{sgn}(p_k)\right]}
\]

Each pole contributes a mode $e^{p_k t}$ in time domain. The contour integral selects which modes contribute based on pole locations relative to $C$.

\subsubsection{Case 2: Gaussian Form}\label{app:case-gaussian}

Consider $H_a(\omega) = e^{-a\omega^2}$ where $a > 0$. Since $H_a(-\omega) = H_a(\omega)$ (even function):
\begin{align*}
I &= \frac{1}{2}\oint_C e^{-a\omega^2} \left[(S_e - izS_o) + (S_e + izS_o)\right] d\omega = S_e \oint_C e^{-a\omega^2} \, d\omega
\end{align*}

The $S_o$ term completely cancels. For a contour along the real axis: $I = S_e \sqrt{\frac{\pi}{a}}$. The Gaussian is the generating function for Hermite polynomials, with time-domain impulse response $h(t) = \frac{1}{\sqrt{4\pi at}} e^{-t^2/(4a)}$ (fundamental solution to the heat equation).

\subsubsection{Case 3: Resonant Functions}\label{app:case-resonant}

Consider $H_a(\omega) = \frac{1}{\omega^2 + \omega_0^2}$. Since this is even, the $S_o$ term vanishes: $I = S_e \oint_C \frac{d\omega}{\omega^2 + \omega_0^2}$. The function has simple poles at $\omega = \pm i\omega_0$ with residues $\pm \frac{1}{2i\omega_0}$. If $C$ encloses only the upper half-plane pole:
\[
I = \frac{\pi S_e}{\omega_0}
\]

This represents a simple harmonic oscillator with time-domain response $h(t) = \frac{\sin(\omega_0 t)}{\omega_0} \Theta(t)$.

\subsubsection{Case 4: Self-Adjoint (Real on Imaginary Axis)}\label{app:case-self-adjoint}

If $H_a(i\theta) \in \mathbb{R}$ for all real $\theta$, then $H_a(-\omega) = \overline{H_a(\omega)}$ (Hermitian symmetry). Writing $H_a(\omega) = H_R(\omega) + iH_I(\omega)$ with $H_R$ even and $H_I$ odd:
\[
I = \oint_C \left[S_e H_R(\omega) + z S_o H_I(\omega)\right] d\omega + i\oint_C \left[S_e H_I(\omega) - z S_o H_R(\omega)\right] d\omega
\]

If $H_a(i\theta)$ is purely real on the imaginary axis ($H_I(i\theta) = 0$) and the contour lies there, the odd component decouples completely: $I = \oint_C S_e H_R(\omega) \, d\omega$. This applies to all causal real-valued systems satisfying $H(s^*) = H^*(s)$, including typical low-pass, band-pass, and physically realizable linear systems.


\section{Shannon Wavelet Derivations}\label{app:shannon-wavelets}

\paragraph{Phase error for complex $z$.} For $z = x + iy$ with $y \neq 0$:
\begin{align*}
e^{iz} &= e^{i(x+iy)} = e^{ix - y}, \\
\Gamma(2N+2, iz) &= \Gamma(2N+2) - \gamma(2N+2, iz),
\end{align*}
where $\gamma(s, w) = \int_0^w t^{s-1} e^{-t} \, dt$ is the lower incomplete gamma function. The imaginary part captures both:
\begin{itemize}
\item \emph{Magnitude error}: from truncating the Taylor series to $N$ terms.
\item \emph{Phase error}: from the complex exponential $e^{iz}$ interacting with the incomplete gamma function.
\end{itemize}

Thus, for applications requiring phase accuracy (e.g., complex filters, IQ modulation), the full complex projector must be used, and the hypergeometric form is \emph{only} a computational simplification for real-valued signals.

\subsection{Frequency Response Contribution $\Phi_n(\omega)$}

The term-specific frequency contribution $\Phi_n(\omega)$ in \eqref{eq:composite_kernel} arises from the Fourier transform of the $n$-th basis monomial $x^{2n}$.

For a polynomial projector:
\[
\widetilde{\mathbb{P}}_N[f](x) = \sum_{n=0}^{N} \widetilde{c}_n x^{2n},
\]
the frequency response is:
\[
\widehat{\widetilde{\mathbb{P}}_N[f]}(\omega) = \sum_{n=0}^{N} \widetilde{c}_n \cdot \mathcal{F}\{x^{2n}\}(\omega),
\]
where $\mathcal{F}\{x^{2n}\}(\omega) = i^{2n} \cdot 2\pi \delta^{(2n)}(\omega)$ is the $(2n)$-th distributional derivative of the Dirac delta.

For practical filter design, we regularize by restricting to a compact domain $[0, x_{\max}]$:
\[
\Phi_n(\omega) = \int_0^{x_{\max}} x^{2n} e^{-i\omega x} \, dx = \frac{\partial^{2n}}{\partial (i\omega)^{2n}} \left[\frac{e^{-i\omega x_{\max}} - 1}{-i\omega}\right].
\]

Using integration by parts $2n$ times:
\[
\Phi_n(\omega) = \sum_{k=0}^{2n} \binom{2n}{k} \frac{(2n-k)!}{\omega^{2n-k+1}} \left[(-i)^k e^{-i\omega x_{\max}} x_{\max}^k - (-i)^k \delta_{k,0}\right],
\]
which exhibits $O(\omega^{-(2n+1)})$ decay, ensuring well-defined frequency localization.

\subsection{Greedy Bit Allocation with Truncation-Quantization Trade-off}

The complete optimization problem is:
\begin{equation}
\min_{\substack{N \in \mathbb{N}_0 \\ \mathbf{m} \in \mathbb{N}_0^{N+1}}} \left[\mathcal{E}_{\text{trunc}}(N) + \mathcal{E}_{\text{quant}}(N, \mathbf{m})\right] \quad \text{s.t.} \quad \sum_{n=0}^{N} m_n \leq B_{\text{total}}.
\tag{P}
\label{eq:full_optimization}
\end{equation}

where:
\begin{align*}
\mathcal{E}_{\text{trunc}}(N) &= \left|\int_0^{x_{\max}} R_N(z) \overline{R(z)} \, dz\right|, \\
\mathcal{E}_{\text{quant}}(N, \mathbf{m}) &= \sum_{n=0}^{N} \frac{S_n}{2^{m_n+1}}.
\end{align*}

\paragraph{Algorithm.}
\begin{enumerate}
\item \textbf{Outer loop (choose $N$):} For $N = 0, 1, 2, \ldots, N_{\max}$:
\begin{enumerate}
\item Compute truncation error $\mathcal{E}_{\text{trunc}}(N)$.
\item Compute sensitivities $S_n$ for $n = 0, \ldots, N$.
\item Solve the inner bit allocation problem for this $N$.
\end{enumerate}
\item \textbf{Inner loop (greedy bit allocation for fixed $N$):}
\begin{enumerate}
\item Initialize $m_n = 1$ for all $n \in \{0, \ldots, N\}$.
\item Set $B_{\text{remaining}} = B_{\text{total}} - (N+1)$.
\item While $B_{\text{remaining}} > 0$:
\begin{itemize}
\item Compute marginal error reduction:
\[
\Delta_n = S_n \left(\frac{1}{2^{m_n+1}} - \frac{1}{2^{m_n+2}}\right) = \frac{S_n}{2^{m_n+2}}.
\]
\item Find $n^* = \arg\max_n \Delta_n$.
\item Update: $m_{n^*} \gets m_{n^*} + 1$, $B_{\text{remaining}} \gets B_{\text{remaining}} - 1$.
\end{itemize}
\item Return $\mathcal{E}_{\text{quant}}(N, \mathbf{m})$.
\end{enumerate}
\item \textbf{Select optimal $N^*$:}
\[
N^* = \arg\min_{N} \left[\mathcal{E}_{\text{trunc}}(N) + \mathcal{E}_{\text{quant}}(N, \mathbf{m}^*(N))\right].
\]
\end{enumerate}

\paragraph{Complexity.} The outer loop runs $O(N_{\max})$ times. For each $N$, the inner greedy algorithm runs $O(B_{\text{total}} \cdot N)$ iterations. Each iteration requires $O(N)$ comparisons to find $\arg\max$. Total: $O(N_{\max} \cdot B_{\text{total}} \cdot N^2)$, which can be optimized to $O(N_{\max} \cdot B_{\text{total}} \cdot N \log N)$ using a priority queue.

\subsection{Residue-Based Decomposition for Hardware Kernels}

Consider a rational kernel:
\[
H(\omega) = \frac{P(\omega)}{Q(\omega)}, \quad Q(\omega) = \prod_{p=1}^{M} (\omega - \omega_p)^{r_p},
\]
where $\omega_p \in \mathbb{C}$ are the poles with multiplicities $r_p$.

The wavelet-projected integral:
\[
I_j(t) = \int_{2^j\pi}^{2^{j+1}\pi} H(\omega) \widetilde{\Phi}_N^{(\mathbf{m})}((i\omega \cdot 2^{-j}t)^2) e^{i\omega t} \, d\omega
\]
can be evaluated by:
\begin{enumerate}
\item Extend the integration domain to a contour $\Gamma_j$ enclosing $[2^j\pi, 2^{j+1}\pi]$ and all poles $\omega_p$ in this range.
\item Apply the residue theorem:
\[
I_j(t) = 2\pi i \sum_{\omega_p \in [2^j\pi, 2^{j+1}\pi]} \operatorname{Res}_{\omega=\omega_p} \left[H(\omega) \widetilde{\Phi}_N^{(\mathbf{m})}((i\omega \cdot 2^{-j}t)^2) e^{i\omega t}\right].
\]t
\end{enumerate}

For a simple pole $\omega_p$ (i.e., $r_p = 1$):
\[
\operatorname{Res}_{\omega=\omega_p} \left[\frac{P(\omega)}{Q(\omega)} F(\omega)\right] = \frac{P(\omega_p)}{Q'(\omega_p)} F(\omega_p),
\]
where $F(\omega) = \widetilde{\Phi}_N^{(\mathbf{m})}((i\omega \cdot 2^{-j}t)^2) e^{i\omega t}$.

Expanding $\widetilde{\Phi}_N^{(\mathbf{m})}(w) = \sum_{n=0}^{N} \frac{a_n}{2^{m_n}} \frac{w^n}{(2n)!}$:
\begin{align*}
F(\omega_p) &= \sum_{n=0}^{N} \frac{a_n}{2^{m_n} (2n)!} \left[(i\omega_p \cdot 2^{-j}t)^2\right]^n e^{i\omega_p t} \\
&= e^{i\omega_p t} \sum_{n=0}^{N} \frac{a_n (-1)^n}{2^{m_n} (2n)!} \left(\frac{\omega_p t}{2^j}\right)^{2n}.
\end{align*}

Thus:
\[
I_j(t) = 2\pi i \sum_{\omega_p} \underbrace{\frac{P(\omega_p)}{Q'(\omega_p)}}_{\text{constant}} \cdot e^{i\omega_p t} \cdot \underbrace{\sum_{n=0}^{N} \frac{a_n (-1)^n}{2^{m_n} (2n)!} \left(\frac{\omega_p t}{2^j}\right)^{2n}}_{\text{dyadic polynomial in } t}.
\]

\subsection{Signal Synthesis Capabilities of Shannon Wavelet Decomposition}

\subsubsection{What Can Be Synthesized}

The Shannon wavelet framework with $n$ scales starting at resolution $j_0$ can exactly represent any bandlimited signal with Fourier support in $[-2^{j_0+n}\pi, 2^{j_0+n}\pi]$. This is the Shannon-Nyquist sampling theorem extended to wavelet bases~\cite{Daubechies1992}.

\paragraph{Fundamental Limitation: Bandlimited Signals Only.}
Unlike time-domain polynomial approximation or windowed sinc functions, the Shannon wavelet \emph{cannot} represent:
\begin{itemize}
\item Sharp transients (step functions, clicks, impulses)
\item Signals with infinite bandwidth
\item Non-smooth waveforms (square waves, sawtooth without filtering)
\end{itemize}

However, \emph{any practical audio or RF signal} is effectively bandlimited by:
\begin{itemize}
\item Physical limits (speaker/microphone bandwidth, antenna bandwidth)
\item Anti-aliasing filters before sampling
\item Nyquist limit of the sampling rate
\end{itemize}

\subsubsection{Concrete Synthesis Examples}

\paragraph{Example 1: Additive synthesis (organ, FM synthesis).}
Synthesize a tone with 8 harmonics at frequency $f_0$:
\[
s(t) = \sum_{k=1}^{8} A_k \sin(2\pi k f_0 t + \phi_k).
\]

For $f_0 = 440$ Hz (A4) with 8 harmonics, highest frequency is $8 \times 440 = 3520$ Hz.

\emph{Shannon wavelet representation:}
\begin{itemize}
\item Choose $j_0$ such that $2^{j_0}\pi / T_s \approx 220$ Hz (below fundamental).
\item 3-scale decomposition covers up to $8 \times 2^{j_0}\pi / T_s \approx 7$ kHz.
\item Each harmonic $k$ lives in scale $j = j_0 + \lfloor \log_2(k) \rfloor$.
\item Total coefficients: $\sim$16 per harmonic $\times$ 8 harmonics = 128 coefficients.
\end{itemize}

\emph{Computational savings vs. time-domain:}
\begin{itemize}
\item Time-domain sinc interpolation: requires $O(L)$ samples over support $L \approx$ several periods.
\item Shannon wavelet: $O(8)$ wavelets (one per harmonic) with compact frequency support.
\item \textbf{Advantage:} Adding/removing harmonics is trivial (set coefficient to zero), enabling real-time additive synthesis with frequency-specific envelope control.
\end{itemize}

\paragraph{Example 2: Subtractive synthesis (sawtooth $\to$ lowpass filter).}
Start with a bandlimited sawtooth (infinite harmonics, amplitude $\propto 1/k$):
\[
\text{sawtooth}(t) = \sum_{k=1}^{\infty} \frac{(-1)^{k+1}}{k} \sin(2\pi k f_0 t).
\]

Bandlimit to 32 harmonics (typical for audio):
\begin{itemize}
\item 5-scale Shannon decomposition ($j_0 = 0$ to $j_0 + 4$) covers 32 octaves.
\item Apply a 2-pole lowpass filter $H(\omega) = \frac{1}{(1 + i\omega/\omega_c)^2}$ with cutoff $\omega_c$.
\item In the wavelet domain, this is a \emph{diagonal scaling} of coefficients:
\[
d_{j,k}^{\text{filtered}} = H(2^{j+1/2}\pi) \cdot d_{j,k}^{\text{original}}.
\]
\item \textbf{Result:} Filter is a single multiplication per scale, not a convolution.
\end{itemize}

\emph{Hardware advantage:}
\begin{itemize}
\item Time-domain FIR: $O(N_{\text{taps}})$ MACs per sample (e.g., 128-tap filter = 128 MACs).
\item Shannon wavelet: 5 complex multiplications (one per scale), then reconstruct.
\item For real-time filter sweeps (cutoff $\omega_c(t)$ varies), only the scale weights change—no filter redesign needed.
\end{itemize}

\paragraph{Example 3: Phase vocoder / time-stretching.}
Time-stretch a signal by factor $\alpha$ without changing pitch:
\begin{itemize}
\item Decompose into Shannon wavelets $\{d_{j,k}\}$.
\item Modify temporal positions: $\psi_{j,k}(t) \to \psi_{j,k}(t/\alpha - k)$ (dilation in time).
\item Frequency content (pitch) is unchanged because wavelets are shift-invariant in frequency.
\end{itemize}

\emph{Traditional phase vocoder:}
\begin{itemize}
\item STFT (Short-Time Fourier Transform): window size trades off time vs. frequency resolution.
\item Shannon wavelets: \emph{native} multi-resolution—low frequencies have long windows, high frequencies have short windows automatically.
\end{itemize}

\paragraph{Example 4: Parametric EQ (audio mixing).}
Boost/cut at specific frequency bands (e.g., bass boost at 100 Hz, treble cut at 8 kHz):
\begin{itemize}
\item Shannon wavelet scale $j$ corresponds to octave band $[2^j\pi, 2^{j+1}\pi]$.
\item To boost scale $j$ by $+6$ dB: multiply all $d_{j,k}$ by $2$.
\item To apply a shelf filter: ramp the gain across scales.
\end{itemize}

\emph{Advantage over IIR/FIR EQ:}
\begin{itemize}
\item IIR (biquad cascade): 2nd-order sections with stability concerns, phase distortion.
\item FIR (Parks-McClellan): linear phase but high order ($\sim$100+ taps for sharp transitions).
\item Shannon wavelet: natural octave bands, perfect reconstruction, no phase distortion within each band.
\end{itemize}

\subsubsection{Quantization Error in Synthesis Context}

For a signal with $n$ scales, total synthesis error is:
\[
\mathcal{E}_{\text{total}} = \underbrace{\sum_{k} |c_{j_0,k}|^2 \mathcal{E}_{\text{quant}}^{(j_0)}}_{\text{base scale}} + \underbrace{\sum_{j=j_0}^{j_0+n-1} \sum_{k} |d_{j,k}|^2 \mathcal{E}_{\text{quant}}^{(j)}}_{\text{detail scales}},
\]
where $\mathcal{E}_{\text{quant}}^{(j)} = \sum_{m} S_{\psi,m}^{(j)} / 2^{m_{\psi,m}^{(j)}+1}$ from Section~4.4.6.

\paragraph{Key insight for synthesis:}
\begin{itemize}
\item High frequencies (large $j$): typically have \emph{small} coefficients $|d_{j,k}|$ (natural signal decay).
\item Low frequencies (small $j$): have \emph{large} coefficients but fewer wavelets needed.
\item The bit allocation vector $\mathbf{m}$ can assign \emph{fewer bits to high-frequency scales} where $|d_{j,k}| \ll 1$, saving hardware resources without audible/visible degradation.
\end{itemize}

\paragraph{Practical synthesis example:}
\begin{itemize}
\item 5-scale decomposition for audio (20 Hz – 20 kHz).
\item Scale 0 ($j_0$): 20–40 Hz (bass), $m_n = 16$ bits (high precision for low distortion).
\item Scale 1: 40–80 Hz, $m_n = 14$ bits.
\item Scale 2: 80–160 Hz, $m_n = 12$ bits.
\item Scale 3: 160–320 Hz, $m_n = 10$ bits.
\item Scale 4: 320–640 Hz and above, $m_n = 8$ bits (high frequencies masked by louder low frequencies).
\item \textbf{Total bits:} $16 + 14 + 12 + 10 + 8 = 60$ bits distributed across 5 scales.
\item \textbf{Uniform quantization:} would use $12 \times 5 = 60$ bits but with \emph{worse} perceptual quality (over-quantizes bass, under-quantizes treble relative to perception).
\end{itemize}

\subsubsection{Comparison: Shannon Wavelets vs. Alternatives}

\begin{table}[h]
\centering
\small
\begin{tabular}{|l|c|c|c|c|}
\hline
\textbf{Method} & \textbf{Freq. Selectivity} & \textbf{Time Localization} & \textbf{Hardware Cost} & \textbf{Phase} \\ \hline
Shannon Wavelet & Excellent & Poor & Low (dyadic ops) & Linear \\ \hline
Haar Wavelet & Poor & Excellent & Very Low & Linear \\ \hline
Daubechies Wavelet & Good & Good & Medium & Nonlinear \\ \hline
FFT/STFT & Excellent & Poor & Medium & Complex \\ \hline
FIR Filter & Good & Medium & High ($N$ taps) & Linear \\ \hline
IIR Filter & Good & Good & Low & Nonlinear \\ \hline
Polynomial (dyadic) & N/A & Excellent & Very Low & N/A \\ \hline
\end{tabular}
\caption{Comparison of signal synthesis methods for bandlimited signals.}
\end{table}

\paragraph{When to use Shannon wavelets:}
\begin{itemize}
\item \textbf{Frequency-domain processing} with minimal computational cost (filtering, EQ, pitch shift).
\item \textbf{Multi-rate systems} where different frequency bands are processed at different rates (e.g., low frequencies at 1 kHz, high frequencies at 48 kHz).
\item \textbf{Adaptive quantization} where bit budget must be allocated across frequency bands based on sensitivity.
\item \textbf{Perfect reconstruction} is required (lossless analysis-synthesis).
\end{itemize}

\paragraph{When NOT to use Shannon wavelets:}
\begin{itemize}
\item \textbf{Transient-heavy signals} (percussion, clicks) — use Daubechies or Haar for better time localization.
\item \textbf{Wideband noise} — time-domain polynomial or direct sampling is simpler.
\item \textbf{Ultra-low latency} — Shannon wavelets have infinite time support (though truncated in practice).
\end{itemize}

\section{Hardware design choices}\label{app:hardware}

\subsection{Kernel design}

\paragraph{Hardware implementation of projector}
\begin{itemize}
\item Precompute constants $C_p = 2\pi i \frac{P(\omega_p)}{Q'(\omega_p)}$ (stored in BRAM).
\item Oscillator: Generate $e^{i\omega_p t}$ using CORDIC in rotation mode (1 DSP per pole).
\item Polynomial: Evaluate the dyadic polynomial using Horner's method ($N$ DSPs or 1 DSP with $N$-cycle latency).
\item Accumulate: Sum over all poles $p$ (1 DSP).
\end{itemize}

For $M$ poles and $N$ terms: $M(N+1) + 1$ DSP slices (parallel) or $N+2$ DSP slices (time-multiplexed over $M$ poles).

This architecture naturally supports frequency-dependent bit allocation: high-sensitivity poles can use larger $m_n$ for their polynomial coefficients, while low-sensitivity poles use reduced precision, all determined by the sensitivity matrix $\mathbf{S}_{\text{wavelet}}$.

\subsection{Fabric design}

\subsection{Synthesis Quality Metrics}

For the quantized Shannon wavelet synthesis $\widetilde{f}^{(N,M,\mathbf{m},\mathbf{m}')}(t)$ from \eqref{eq:proj_approx}:

\paragraph{Signal-to-Quantization-Noise Ratio (SQNR):}
\[
\text{SQNR} = 10 \log_{10} \frac{\|f\|_{L^2}^2}{\|f - \widetilde{f}^{(N,M,\mathbf{m},\mathbf{m}')}\|_{L^2}^2} \quad \text{(dB)}.
\]

For the greedy bit allocation (Appendix A.3) with total budget $B_{\text{total}} = 60$ bits across 5 scales:
\begin{itemize}
\item Typical SQNR: 70–90 dB (depending on signal).
\item Comparable to 12-bit uniform quantization on the time-domain signal but with \emph{perceptually optimized} distribution.
\end{itemize}

\paragraph{Total Harmonic Distortion (THD):}
For synthesis of harmonic signals (musical tones):
\[
\text{THD} = \frac{\sqrt{\sum_{k=2}^{\infty} P_k}}{P_1} \times 100\%,
\]
where $P_k$ is the power in the $k$-th harmonic.

With $N=4$ polynomial terms per wavelet and $m_n \in [8, 16]$ bits:
\begin{itemize}
\item THD $< 0.01\%$ for fundamental + first 8 harmonics (professional audio quality).
\item THD $< 0.1\%$ for 32 harmonics (CD quality).
\end{itemize}

\paragraph{Bandwidth efficiency:}
\begin{itemize}
\item Shannon wavelets are \emph{optimal} for bandlimited signals in the sense that the number of coefficients $\{c_{j_0,k}, d_{j,k}\}$ equals the number of degrees of freedom (Shannon's sampling theorem).
\item For a signal bandlimited to $B$ Hz sampled at $f_s = 2B$, the Shannon wavelet decomposition requires $\sim f_s$ coefficients per second, which is \emph{information-theoretically optimal}.
\end{itemize}

\subsection{Hardware Resource Estimates for Shannon Wavelet Implementation}

\paragraph{FPGA resource model.} Consider a Xilinx Virtex-7 FPGA with the following primitives:
\begin{itemize}
\item DSP48E1 slices: $18 \times 25$-bit multipliers, 48-bit accumulators.
\item Block RAM (BRAM): 36 Kb dual-port memory blocks.
\item Look-up tables (LUTs): 6-input logic functions.
\end{itemize}

\paragraph{Scaling function computation.} For $\widetilde{\phi}_{j_0,k}^{(N,\mathbf{m})}(t)$:
\[
\widetilde{\phi}_{j_0,k}^{(N,\mathbf{m})}(t) = 2^{j_0/2} \sum_{n=0}^{N} \frac{a_n}{2^{m_n}} \left[\pi(2^{j_0}t - k)\right]^{2n}.
\]

\emph{Resources per basis function:}
\begin{itemize}
\item $(N+1)$ coefficient storage: $(N+1) \times \max_n(m_n + \lceil \log_2 |a_n| \rceil)$ bits in BRAM.
\item Power computation $u^{2n}$ for $u = \pi(2^{j_0}t - k)$: $\lceil \log_2(N+1) \rceil$ DSP slices (using Horner's method or parallel squaring tree).
\item Multiply-accumulate: $(N+1)$ DSP slices or time-multiplexed over $\lceil (N+1)/P \rceil$ cycles for $P$ parallel DSPs.
\item Dyadic scaling $\frac{a_n}{2^{m_n}}$: implemented as barrel shifter using $O(m_n)$ LUTs (no DSP needed).
\end{itemize}

\emph{Example (N=4, $m_n \in [8, 16]$):}
\begin{itemize}
\item Coefficients: $5 \times 24$ bits $= 120$ bits $\approx 0.003$ BRAM.
\item Power computation: 3 DSP slices (parallel).
\item MAC: 5 DSP slices (fully parallel) or 1 DSP slice (5-cycle latency).
\item Total per basis: 8 DSP, 120-bit BRAM, $\sim$50 LUTs.
\end{itemize}

For a 3-scale wavelet decomposition ($j_0, j_0+1, j_0+2$) with $K=32$ spatial positions per scale:
\begin{itemize}
\item Total DSP: $8 \times 32 \times 3 = 768$ DSPs (parallel) or $8 \times 3 = 24$ DSPs (time-multiplexed over 32 cycles).
\item Total BRAM: $120 \times 32 \times 3 = 11520$ bits $\approx 0.32$ BRAM blocks.
\item Throughput: 1 sample per cycle (parallel) or 1 sample per 32 cycles (multiplexed).
\end{itemize}

Virtex-7 XC7VX485T has 2800 DSP slices and 1030 BRAM blocks, so even the fully parallel implementation uses $\sim$27\% DSP and $<$1\% BRAM.

\section{Code snippets or heuristics}

\subsection{'compact' sinc implementation}

\begin{python}
def compact_sinc(x: np.ndarray, dtype: np.dtype) -> np.ndarray[np.floating]:
    x = np.asarray(x, dtype=dtype)
    result = np.zeros_like(x)
    x_nz = x[x != 0]
    if x_nz.size > 0:
        abs_x_nz = np.abs(x_nz)
        lobe_number = np.floor(abs_x_nz)
        lobe_index = abs_x_nz - lobe_number
        envelope = 1.0 / (np.pi * np.abs(x_nz))
        sign = (-1) ** lobe_number
        canonical_values = _canonical_sinc_lobe_lookup(lobe_index)
        result[x != 0] = sign * envelope * canonical_values
    result[x == 0] = 1.0
    return result

def _canonical_sinc_lobe_lookup(lobe_index: np.ndarray) -> np.ndarray[np.floating]:
    # Convert lobe_index from [0, 2) to [0, pi) for sinc calculation
    t = lobe_index * np.pi
    result = np.where(t == 0, 1.0, np.sin(t))
    return result
\end{python}

\subsection{comparison of cos power series in c}\label{app:cos-comparison}

\begin{clang}
static inline double cse_Polynomial_Cos_Deg6(float x) {
  double xd = (double)x;
  double z = xd * xd;
  double t = fma(-1.3888888888888889e-3, z, 4.1666666666666664e-2);
  t = fma(t, z, -5.0e-1);
  t = fma(t, z, 1.0);
  return t;
}


static inline double cse_Polynomial_Cos_Deg12(float x) {
  double xd = (double)x;
  double z = xd * xd;
  double t = fma(2.08767569878681e-09, z, -2.7557319223985888e-07);
  t = fma(t, z, 2.4801587301587302e-05);
  t = fma(t, z, -1.3888888888888889e-03);
  t = fma(t, z, 4.1666666666666664e-02);
  t = fma(t, z, -5.0e-1);
  t = fma(t, z, 1.0);
  return t;
}


static inline double cse_Dyadic_Polynomial_Cos_Deg12(float x) {
  double xd = (double)x;
  double z = xd * xd;
  double t = fma(-2.0563602447509766e-05, z, -1.3113021850585938e-06);
  t = fma(t, z, 2.3901462554931641e-05);
  t = fma(t, z, -0.0013893246650695801);
  t = fma(t, z, 0.041666388511657715);
  t = fma(t, z, -0.50000017881393433);
  t = fma(t, z, 0.99999988079071045);
  return t;
}
\end{clang}

\subsection{introducing concurrency into power series}\label{app:simd-vectorization}

\begin{python}
  def _convert_lines_to_simd(self, lines: list[str], target: SIMD = SIMD.AVX512) -> list[str]:
        """Annotate per-stage lines with AVX-512 FMAs packed by common operand.

        Looks for lines of form 'lhs = fma(a, b, c)' and groups those sharing the
        same 'b' into packs of up to 8 lanes (double). Emits a single fmadd
        comment per pack, then preserves all original scalar lines and order.
        """
        if target is not SIMD.AVX512:
            raise NotImplementedError('Only AVX512 is implemented')

        assign_pat = re.compile(r'^\s*([A-Za-z_]\w*)\s*=\s*(.+)$')
        fma_pat = re.compile(r'^fma\(\s*(?P<a>.+?)\s*,\s*(?P<b>.+?)\s*,\s*(?P<c>.+?)\s*\)\s*$')

        # Parse entries and track original indices
        entries: list[tuple[int, str, str, str, str]] = []  # (idx, lhs, a, b, c)
        for idx, ln in enumerate(lines):
            m = assign_pat.match(ln)
            if not m:
                continue
            lhs, rhs = m.group(1), m.group(2).strip()
            fm = fma_pat.match(rhs)
            if fm:
                a = fm.group('a').strip()
                b = fm.group('b').strip()
                c = fm.group('c').strip()
                entries.append((idx, lhs, a, b, c))

        # Group by common b operand
        groups: dict[str, list[tuple[int, str, str, str]]] = {}
        for idx, lhs, a, b, c in entries:
            groups.setdefault(b, []).append((idx, lhs, a, c))

        # Build code to insert at the earliest line index of each group
        inserts_at: dict[int, list[str]] = {}
        drop_indices: set[int] = set()

        def chunk(seq, n):
            for i in range(0, len(seq), n):
                yield seq[i:i + n]

        def pad_eight(values: list[str]) -> list[str]:
            v = list(values)
            while len(v) < 8:
                v.append('0.0')
            return v[:8]

    # Use a global, monotonically-increasing id to ensure unique symbol names across calls
        for common, group in groups.items():
            if len(group) < 2:
                continue  # skip singletons

            # Sort by original order for stable packing and mark to drop
            group_sorted = sorted(group, key=lambda t: t[0])
            first_idx = group_sorted[0][0]
            lines_for_group: list[str] = []

            for ch in chunk(group_sorted, 8):
                ls = [lhs for (idx, lhs, _, _) in ch]
                as_ = pad_eight([a for (_, _, a, _) in ch])
                cs_ = pad_eight([c for (_, _, _, c) in ch])
                lanes = ', '.join(ls)
                lines_for_group.append(f"// avx512 fmadd pack: common={common} -> lanes {lanes}")
                pack_id = self._simd_gencounter
                self._simd_gencounter += 1
                lines_for_group.append(f"__m512d __simd_va{pack_id} = _mm512_setr_pd({', '.join(as_)});")
                lines_for_group.append(f"__m512d __simd_vc{pack_id} = _mm512_setr_pd({', '.join(cs_)});")
                lines_for_group.append(f"__m512d __simd_vx{pack_id} = _mm512_set1_pd({common});")
                lines_for_group.append(f"__m512d __simd_vout{pack_id} = _mm512_fmadd_pd(__simd_va{pack_id}, __simd_vx{pack_id}, __simd_vc{pack_id});")
                lines_for_group.append(f"double __simd_pack{pack_id}[8];")
                lines_for_group.append(f"_mm512_storeu_pd(__simd_pack{pack_id}, __simd_vout{pack_id});")
                for lane, lhs in enumerate(ls):
                    lines_for_group.append(f"{lhs} = __simd_pack{pack_id}[{lane}];")

            inserts_at[first_idx] = lines_for_group
            # Mark all indices in this group to drop (they are replaced)
            for idx, *_ in group:
                drop_indices.add(idx)

        # Emit original lines with inserts placed at calculated positions; drop replaced scalars
        out: list[str] = []
        for idx, ln in enumerate(lines):
            if idx in inserts_at:
                out.extend(inserts_at[idx])
            if idx in drop_indices:
                continue
            out.append(ln)
        return out
\end{python}\label{eq:simd_power_series}
